\section{Background \label{sec:background}}

Software artifacts are naturally relational and multi-typed. A system contains folders, source files or classes, and finer-grained program elements such as methods and fields, connected through relations that include folder containment, file-to-member composition, and static code-level dependencies such as calls, inheritance, type usage, and field access.
 
To preserve these distinctions, we represent each software system as a heterogeneous software graph $G=(V, E,\mathcal{A},\mathcal{R})$. Here, $V$ denotes the set of software artifact nodes and $E$ denotes the set of relations between them. Each node $v \in V$ is assigned a node type through the mapping function $\tau(v): V \rightarrow \mathcal{A}$, where $\mathcal{A}$ contains folder, file, and aggregate member types in our most detailed configuration. Similarly, each edge $e \in E$ is assigned a relation type through $\phi(e): E \rightarrow \mathcal{R}$, capturing folder hierarchy, the association between folders and classes, the association between classes and finer-grained program elements, and static code dependencies such as method calls, inheritance, and imports.
 
In this graph, the meaning of a relation is not characterized by the relation type alone. It also depends on the types of the source and target nodes that the edge connects. Therefore, for an edge $e_{s,t}$ from source node $s$ to target node $t$, we describe the typed interaction as a meta-relation $\langle \tau(s), \phi(e_{s,t}), \tau(t) \rangle$. For example, a $\langle \text{folder},\ contains,\ \text{file} \rangle$ relation captures project organization, whereas a $\langle \text{file},\ Import,\ \text{file} \rangle$ relation captures code-level coupling. Thus, the same graph can preserve both structural organization and implementation dependencies without treating them as identical relations. A homogeneous graph, with one node type and one relation type, discards this distinction, and RQ3 uses it to test whether the typed, directed form matters.
 
HGT~\cite{hgt} exploits such meta-relations during message passing by conditioning attention and message transformation on the source node type, edge type, and target node type. GAT~\cite{gat} also uses attention-based message passing, but its standard formulation does not explicitly distinguish node or relation types. HGT can also operate on a graph with a single node type and a single relation type, although this does not make its attention mechanism equivalent to that of GAT, and RQ3 compares both on such a graph. The initial node features determine what information is available to these models, so graph structure and feature representation describe different parts of a mapping model, even when both are derived from dependencies.
 
A lexical embedding trained on identifier tokens captures naming conventions but not a file's position within the dependency structure or the syntactic and behavioral content of its implementation. An embedding trained directly on the dependency graph, such as Node2Vec~\cite{node2vec}, instead encodes a file's structural neighborhood without reference to its content, and such embeddings have been shown to capture dependency-based cohesion effectively in architecture recovery~\cite{negar,structsar}.
 
Pretrained code encoders such as UniXcoder~\cite{unixcoder} address the content gap left by lexical features by learning general-purpose representations of source code from large corpora, but their pretraining objective is not specific to architectural grouping: two files can be lexically and syntactically similar while belonging to different modules. Contrastive learning offers a way to adapt such an encoder to a task-specific notion of similarity by training the model to place designated positive pairs closer in the embedding space than unrelated in-batch examples, typically through an InfoNCE objective~\cite{infonce}. Using positive pairs derived from dependencies encourages the embedding space to reflect dependency proximity. It does not guarantee that this proximity coincides with the reference modules.